\chapter{第4章\texorpdfstring{\hspace{0.5em}}{ }流体力学简介}

气体（gas）和液体（liquid）统称为\textbf{流体}（fluid），它们属于容易形变的连续分布的介质。流体的基本特征是具有\textbf{流动性}，即各部分之间极容易发生相对运动：气体易被压缩、弥漫于整个容器空间、不存在自由表面；液体不易被压缩、有一定体积、能形成自由表面。流体力学是连续介质力学的一部分，以牛顿运动定律为基础研究流体运动的宏观规律。本章提出理想流体模型，讨论流体运动的描述、连续性方程与伯努利方程，并简要介绍黏性流体的运动规律。

\section{流体运动的描述}
    \subsection{速度场和定常流动}
    由于流体的流动性，各部分质元的运动不一定相同。因此处理流体运动时，并不着眼于各个流体质元在各时刻的运动状态，而是考察流体所占据的空间中各点，研究流体各质元流经这些点时所具有的速度、密度、压强等，以及这些量随时间变化的规律（欧拉法）。

    流体运动过程中每一瞬时，流体所占据的空间每一点都具有一定的流速，这种流速的空间分布称为\textbf{速度场}（velocity field），用矢量函数 $\boldsymbol{v}(x,y,z,t)$ 描述：
    \[
    \boldsymbol{v}=\boldsymbol{v}(x,y,z,t)
    \]
    它表示 $t$ 时刻空间任一点 $(x,y,z)$ 处流体质元的速度。同理还有密度场 $\rho(x,y,z,t)$、压强场 $p(x,y,z,t)$、温度场 $T(x,y,z,t)$ 等。

    若速度场不随时间变化，即
    \keyformula{\frac{\partial\boldsymbol{v}}{\partial t}=0\quad\Longleftrightarrow\quad \boldsymbol{v}=\boldsymbol{v}(x,y,z)}
    这种流动称为\textbf{定常流动}（steady flow，又称恒定流动）。定常流动中，空间任一确定位置的流速、压强、密度均不随时间改变，但不同位置的流速一般不同。
    \corepoint{“定常”是对空间固定点而言的：同一流体质元在流动中完全可以加速（因为它流到了流速更大的位置），这与“流速处处相同”不是一回事。}

    \subsection{流线和流管}
    为了形象地描述流体运动，在流体中作一系列曲线，使曲线上每一点的切线方向与该处流体质元的速度方向一致，这些曲线称为\textbf{流线}（streamline）。
    \begin{enumerate}
        \item 某一时刻每个质元只有一个确定的速度方向，所以流线\textbf{不会相交}；
        \item 流速大的地方流线密，流速小的地方流线疏；
        \item 定常流动中流线的形状稳定不变，且与流体质元运动的轨迹（迹线）重合。
    \end{enumerate}

    在流体中由流线围成的管状空间称为\textbf{流管}（tube of flow）。由于流线不相交，流管内外的流体都不会穿越管壁流动：管内流体不能流到管外，管外流体也不能流入管内。对定常流动，流线与流管的形状都不随时间变化，流管的行为与真实的固体管道相似，这是把流体问题化为“管道流动”处理的基础。

\section{连续性方程和伯努利方程}
    \subsection{理想流体}
    液体在外力作用下体积变化很小，一般可忽略其可压缩性；气体虽然容易被压缩，但流动性强，在定常情况下密度一般不会明显变化，故对低速气体（流速远小于其中声速）也可不考虑其压缩性。

    当相邻两部分流体有相对运动时，沿接触面的切向会产生内摩擦力，又称流体的\textbf{黏性力}（viscous force）。对黏性很弱、流程又不长的流体，内摩擦力引起的机械能损耗可以忽略，即可认为流体没有黏性，各流层间只存在法向相互作用力而无切向相互作用力。

    为了简化问题、突出流体的流动性，研究中用\textbf{理想流体}（ideal fluid）模型替代真实流体：
    \begin{center}
    \begin{tabular}{ll}
    \toprule
    理想流体 & 绝对不可压缩、完全没有黏性的流体 \\
    \midrule
    真实流体 & 既具有流动性，又具有可压缩性和黏性 \\
    \bottomrule
    \end{tabular}
    \end{center}

    \subsection{理想流体的连续性方程}
    在作定常流动的理想流体中取一细流管，流管很细，可认为同一截面上各点流速相同。任取两个垂直于流管的截面 $a$、$b$，面积分别为 $\Delta S_{a}$、$\Delta S_{b}$，流速分别为 $v_{a}$、$v_{b}$。由于理想流体不可压缩，单位时间内流入 $a$ 截面的流体体积必等于从 $b$ 截面流出的体积：
    \keyformula{v_{a}\Delta S_{a}=v_{b}\Delta S_{b}\quad\Longrightarrow\quad v\Delta S=\text{常量}}
    由于 $a$、$b$ 是任意选取的，故对同一流管的任一横截面都成立。式中 $v\Delta S$ 是单位时间内通过该截面的流体体积，称为\textbf{体积流量}（简称流量），记作
    \keyformula{Q_{V}=v\Delta S=\text{常量}\quad\left[\mathrm{m^{3}\cdot s^{-1}}\right]}
    在定常流动下各点密度不随时间变化，故\textbf{质量流量}也守恒：
    \[
    Q_{m}=\rho v\Delta S=\text{常量}
    \]
    式 $v\Delta S=$ 常量称为\textbf{流体的连续性方程}（continuity equation of fluid），它是质量守恒定律在流体力学中的具体表述。若流管共包围 $n$ 根流线，则截面处的流线密度为 $n/\Delta S$，故有
    \[
    \frac{n}{\Delta S_{a}}:\frac{n}{\Delta S_{b}}=v_{a}:v_{b}
    \]
    即\textbf{流线密度与流速成正比}——这正是“流速大处流线密”的定量表述。

    \corepoint{连续性方程的结论是“截面大处流速小、截面小处流速大”。水管口用手指挡住一部分后水射得更远，正是截面变小、流速变大所致。}

    \subsection{理想流体的伯努利方程}
    伯努利方程（Bernoulli equation）是理想流体定常流动的动力学方程，表明沿一条流线各点的压强、流速和高度三者关系的基本规律，1738 年由瑞士物理学家伯努利建立。

    \textbf{推导（功能原理）}：在重力场中作定常流动的理想流体中取一细流管，取 $a$、$b$ 两截面之间的一段流体为研究对象，经时间 $\Delta t$ 后这段流体流到 $a'b'$ 位置。由于不可压缩，$aa'$ 段与 $bb'$ 段的体积相等、质量相等（均为 $\Delta m=\rho\Delta V$），故机械能的变化就是这两段流体机械能之差：
    \[
    \Delta E=\left(\frac{1}{2}\Delta mv_{2}^{2}+\Delta mgh_{2}\right)
    -\left(\frac{1}{2}\Delta mv_{1}^{2}+\Delta mgh_{1}\right)
    \]
    理想流体无黏性，流管外流体对这段流体只有垂直于流管壁的压力，不做功；只有两端截面上的压力做功：
    \[
    W=p_{1}\Delta S_{1}\Delta l_{1}-p_{2}\Delta S_{2}\Delta l_{2}
    =(p_{1}-p_{2})\Delta V
    \]
    由功能原理 $W=\Delta E$，并注意 $\Delta m=\rho\Delta V$，得
    \keyformula{p_{1}+\frac{1}{2}\rho v_{1}^{2}+\rho gh_{1}
    =p_{2}+\frac{1}{2}\rho v_{2}^{2}+\rho gh_{2}}
    因 $a$、$b$ 位置任意，对同一流管（细化为一条流线）的任一截面有
    \keyformula{p+\frac{1}{2}\rho v^{2}+\rho gh=\text{常量}}
    这就是\textbf{伯努利方程}。工程上常除以 $\rho g$ 写成“水头”形式：
    \[
    \frac{p}{\rho g}+h+\frac{v^{2}}{2g}=\text{常量}
    \]
    其中 $\dfrac{v^{2}}{2g}$、$\dfrac{p}{\rho g}$、$h$ 分别称为\textbf{速度水头}（速头）、\textbf{压力水头}（压头）和\textbf{位置水头}（静头），三者之和称为总水头。

    \textbf{适用条件与要点}：
    \begin{enumerate}
        \item 理想流体、定常流动、不可压缩、沿同一细流管（同一条流线）成立；
        \item 未计入其他能量损失；
        \item 式本质是\textbf{单位体积流体的机械能守恒}：$p$ 与单位体积的压强能对应，$\frac{1}{2}\rho v^{2}$ 是单位体积动能，$\rho gh$ 是单位体积重力势能；
        \item $h_{1}$、$h_{2}$ 必须相对\textbf{同一参考平面}度量，两参考点应在同一流线上；
        \item 取 $v=0$ 时即退化为流体静力学结论 $p+\rho gh=$ 常量，故流体静力学是流体动力学的特例。
    \end{enumerate}
    \corepoint{由伯努利方程可直接得到“等高流动时流速大处压强小”：$h$ 相同时 $p+\frac{1}{2}\rho v^{2}=$ 常量。这是喷雾器、汽油雾化器、水流抽气机、飞机升力等现象的共同原理。}

        \subsubsection{应用模型一：小孔流速（托里拆利）}
        大容器侧壁下方开一直径为 $d$ 的小孔，液面距孔高 $h$。取从液面到小孔的一条流线，液面处 $v_{1}\approx0$（容器截面远大于孔），液面与小孔处压强均为大气压 $p_{0}$，以小孔所在水平面为参考面，由伯努利方程
        \[
        p_{0}+\rho gh=p_{0}+\frac{1}{2}\rho v_{2}^{2}
        \]
        得小孔处流速与小孔流量
        \keyformula{v=\sqrt{2gh},\qquad Q_{V}=\frac{\pi d^{2}}{4}\sqrt{2gh}}
        泄流水的动能功率为 $P=\dfrac{1}{2}\rho Q_{V}v^{2}$。

        \subsubsection{应用模型二：文特利流量计}
        文特利（venturi）流量计是两端粗（截面 $S_{1}$）、中间细（截面 $S_{2}$）的短管，粗、细管处接 U 形压强计，串接在水平管道上即可测流量。取管轴线为参考基准线，则 $h_{1}=h_{2}$，由伯努利方程与连续性方程
        \[
        p_{1}+\frac{1}{2}\rho v_{1}^{2}=p_{2}+\frac{1}{2}\rho v_{2}^{2},
        \qquad S_{1}v_{1}=S_{2}v_{2}
        \]
        解得两截面压强差与流量
        \keyformula{p_{1}-p_{2}=\frac{1}{2}\rho\left(v_{2}^{2}-v_{1}^{2}\right),
        \qquad
        Q_{V}=S_{1}S_{2}\sqrt{\frac{2\left(\rho_{\text{汞}}-\rho\right)gh}{\rho\left(S_{1}^{2}-S_{2}^{2}\right)}}}
        其中压强差由 U 形压强计液面高度差 $h$ 给出：$p_{1}-p_{2}=(\rho_{\text{汞}}-\rho)gh$。实际使用时还须按真实流体情况加修正系数。

\section{黏性流体的运动}
    理想流体无黏性，流动中没有机械能损耗。但生产实际中流体的内摩擦力往往不可忽视（如石油在输油管中的运动）；即使是黏性很小的流体，在长距离输送过程中内摩擦引起的机械能损耗也不能忽略。

    \subsection{流体的黏性定律}
    黏性较大的流体在直径不大的管中以较慢的速度流动时，出现各层流速不同的运动状态，各流层之间不相混杂，这种流动称为\textbf{层流}（laminar flow，又称片流）。层流的特点是只有切向速度、没有径向速度。

    设相距 $\Delta y$ 的两流层分别以速度 $v_{1}$、$v_{2}$ 沿 $x$ 轴正方向流动，则流速在 $y$ 方向上的变化率
    \[
    \frac{\Delta v_{x}}{\Delta y}\ \longrightarrow\ \frac{\mathrm{d}v_{x}}{\mathrm{d}y}
    \]
    称为\textbf{速度梯度}，表示由一流层过渡到另一流层时速度变化的快慢。作用在两流层间的内摩擦力沿层面切线方向：速度大的一层受力方向与运动方向相反，速度小的一层受力方向与速度方向一致。在接触面上取面积 $\Delta S$，则内摩擦力大小为
    \keyformula{f=\eta\frac{\mathrm{d}v}{\mathrm{d}y}\Delta S,\qquad
    \tau=\frac{f}{\Delta S}=\eta\frac{\mathrm{d}v}{\mathrm{d}y}}
    称为流体的\textbf{黏性定律}（牛顿黏性定律）。其中 $\eta$ 称为\textbf{黏度}（黏滞系数、内摩擦系数），与流体的性质和温度有关：
    \begin{itemize}
        \item 同一种\textbf{液体}，温度升高，黏度急剧减小；
        \item 同一种\textbf{气体}，温度升高，黏度变大。
    \end{itemize}
    可见液体与气体虽宏观规律相同，微观机理不同。单位：帕斯卡秒（$\mathrm{Pa\cdot s}$），也常用泊（$\mathrm{P}$），$1\ \mathrm{P}=0.1\ \mathrm{Pa\cdot s}$。静止流体中不存在黏性力。

    \subsection{实际流体定常流动的伯努利方程}
    实际流体有黏性，流动中内摩擦力使一部分机械能转变为热能、声能等，因而必须把内摩擦力的功计入。设单位体积流体从 $a$ 位置流到 $b$ 位置时内摩擦力做的功为 $W$（$W<0$），则由功能原理
    \keyformula{p_{1}+\rho gh_{1}+\frac{1}{2}\rho v_{1}^{2}
    =p_{2}+\rho gh_{2}+\frac{1}{2}\rho v_{2}^{2}+w}
    式中 $w$ 为单位体积流体的机械能损失；写成水头形式
    \[
    \frac{p_{1}}{\rho g}+h_{1}+\frac{v_{1}^{2}}{2g}
    =\frac{p_{2}}{\rho g}+h_{2}+\frac{v_{2}^{2}}{2g}+h_{w},
    \qquad h_{w}=\frac{w}{\rho g}
    \]
    $h_{w}$ 称为\textbf{损头}（因内摩擦而损失的水头）。

    对水平圆管（$h_{1}=h_{2}$、$v_{1}=v_{2}$）有
    \keyformula{w=p_{1}-p_{2}}
    即在水平管道中维持黏性流体的定常流动，必须保持一定的压强差。这与理想流体不同：理想流体不需要压强差即可流过均匀水平管，而黏性流体因受管壁黏性力的拖拽，必须由压强差来补偿能量损失。

    \subsection{湍流}
    当流体流速很大或流动横截面的线度增大时，流体除向前流动外还有局部的横向运动，各部分相互扰动，失去层流的性质而变成不规则的运动，这种运动称为\textbf{湍流}（turbulent flow，又称紊流）。其特征是流线混杂紊乱、有垂直于管轴的分速度、出现漩涡。

    英国物理学家雷诺通过实验发现，黏性流体由层流转变为湍流可由一个量纲为一的数——\textbf{雷诺数}（Reynold's number）$Re$ 判断：
    \keyformula{Re=\frac{\rho vl}{\eta}}
    式中 $\rho$ 为流体密度，$\eta$ 为黏度，$v$ 为流体在横截面上的平均速率，$l$ 为物体的特征几何线度（管道直径、河渠深度、机翼厚度、障碍物线度等）。对一定形状的物体（不论大小），层流转为湍流的雷诺数有一个临界值 $Re_{c}$。

    \begin{center}
    \begin{tabular}{lll}
    \toprule
    流动状态 & 圆管中 $Re$ & 特征 \\
    \midrule
    层流 & $Re<2000$ & 各流层互不混掺，只有黏性阻力 \\
    过渡状态 & $2000<Re<4000$ & 不稳定，时层时湍 \\
    湍流 & $Re>4000$ & 流层混掺、有涡旋，除黏性阻力还有惯性阻力 \\
    \bottomrule
    \end{tabular}
    \end{center}
    湍流发生时内摩擦力显著增大、流量减小；由于湍流阻力比层流阻力大得多，工程上应尽量避免或利用湍流（如强化换热）时需具体分析。雷诺数在水工模型实验和风洞实验中十分重要。

    \subsection{斯托克斯定律}
    当物体在黏性流体中运动、速度不大且形状是流线型时，没有涡旋发生，紧靠物体的流体被物体带走，物体受到的阻力是黏性力。斯托克斯从实验中总结出：半径为 $r$ 的圆球以速度 $v$ 在黏度为 $\eta$ 的流体中运动时，受到的阻力为
    \keyformula{f=6\pi\eta rv}
    称为\textbf{斯托克斯定律}（Stoke's law）。对非球形物体，定律的形式仍成立，只是对应的黏度值不同。

    物体在黏性流体中下落时受重力、浮力和黏性力作用，黏性力随速度增大而增大；当浮力与黏性力之和与重力平衡时，下落速度达到最大值，称为\textbf{收尾速度}（终端速度）。密度为 $\rho$、半径为 $r$ 的小球在密度为 $\rho'$、黏度为 $\eta$ 的液体中下落的收尾速度为
    \keyformula{v_{T}=\frac{2\left(\rho-\rho'\right)gr^{2}}{9\eta}}
    该式是沉降法测黏度、颗粒粒径分析（斯托克斯沉降分析）的依据：
    \[
    \eta=\frac{2\left(\rho-\rho'\right)gr^{2}}{9v_{T}},
    \qquad
    r=\sqrt{\frac{9\eta v_{T}}{2\left(\rho-\rho'\right)g}}
    \]

    \subsection{泊肃叶定律}
    为确定黏性流体通过圆管的流量，法国科学家泊肃叶从实验中总结出：黏性流体在水平圆管中作稳定分层流动（层流）时，流量与管两端压强差、管半径的四次方成正比，与管长、黏度成反比。
    \keyformula{Q_{V}=\frac{\pi R^{4}\left(p_{1}-p_{2}\right)}{8\eta l}}
    称为\textbf{泊肃叶定律}（Poiseuille law）。式中 $l$ 是管道长度，$p_{1}-p_{2}$ 是管道两端流体的压强差，$R$ 是管道半径，$\eta$ 是流体黏度。$p_{1}-p_{2}$ 越大、$R$ 越大、$\eta$ 越小、$l$ 越小，公式越准确（即越接近层流）。

    层流时圆管截面上的流速按抛物线分布，管轴处最大、管壁处为零：
    \keyformula{v(r)=\frac{p_{1}-p_{2}}{4\eta l}\left(R^{2}-r^{2}\right),\qquad
    \bar{v}=\frac{Q_{V}}{\pi R^{2}}=\frac{R^{2}\left(p_{1}-p_{2}\right)}{8\eta l}}
    即平均流速等于最大流速 $v_{\max}=\dfrac{R^{2}(p_{1}-p_{2})}{4\eta l}$ 的一半。维持该流动所需的功率为
    \[
    P=\left(p_{1}-p_{2}\right)Q_{V}
    \]
    \corepoint{泊肃叶定律中流量与半径的\textbf{四次方}成正比，这是它最重要的推论：血管半径缩小一半，在同样压强差下流量只有原来的 $1/16$——这正是动脉硬化显著增加心脏负担的原因。}
